\chapter{Note on the Use of Artificial Intelligence}

Within the scope of this thesis, tools based on \gls{acr:AI} were used. Table \ref{tab:anhang_uebersicht_KI_werkzeuge} provides an overview of the tools used and their respective purpose of application.

\begin{table}[hbt]	
	\centering
	\renewcommand{\arraystretch}{1.5}	% Scales the row height of the table
	\captionabove[List of artificial intelligence based tools used]{List of \gls{acr:AI}-based tools used}
	\label{tab:anhang_uebersicht_KI_werkzeuge}
	\begin{tabular}{>{\raggedright\arraybackslash}p{0.3\linewidth} >{\raggedright\arraybackslash}p{0.65\linewidth}}
		\textbf{Tool} & \textbf{Description of Use}\\
		\hline 
		\hline
		DeepL	&	\vspace{-\topsep}
		\begin{itemize}[noitemsep,topsep=0pt,partopsep=0pt,parsep=0pt] 
			\item Translation and drafting assistance (German <--> English)
			\item Translation Abstract chapter
		\end{itemize} \\ 
		Gemini	&	\vspace{-\topsep}
		\begin{itemize}[noitemsep,topsep=0pt,partopsep=0pt,parsep=0pt] 
			\item Assistance with understanding fundamental concepts
			\item Finding Literature Sources
		\end{itemize} \\ 
		Microsoft Copilot 365	&	\vspace{-\topsep}
		\begin{itemize}[noitemsep,topsep=0pt,partopsep=0pt,parsep=0pt] 
			\item Chapter structuring
			\item Drafting assistance
			\item Grammar and spelling checks
			\item Formatting (bullet points, tables, formulas, equations) in LaTeX (overall)
			\item Generation of LaTeX code for Figures \ref{fig:testbed_topology}, \ref{fig:nv1},  \ref{fig:nv2}, \ref{fig:nv3} based on text descriptions
			\item Assistance with understanding fundamental concepts
			\item Code ideas / error analysis for test cases
		\end{itemize} \\ 
		Github Copilot	&	\vspace{-\topsep}
		\begin{itemize}[noitemsep,topsep=0pt,partopsep=0pt,parsep=0pt] 
			\item LaTeX code: Tables
			\item Abbreviations
			\item Formatting (bullet points, tables, formulas, equations) in LaTeX (overall)
		\end{itemize} \\ 
		\hline 
	\end{tabular} 
\end{table}

\clearpage



\chapter{Detailed Requirement Derivation Matrices}
\label{app:requirement_derivation}


This appendix provides the complete, requirement-by-requirement documentation underlying the derivation methodology described in Section \ref{subsec:derivation_mapping}. It contains 14 requirements derived from obligations addressed to machinery and related products under \hyperref[subsec:eu_machinery_regulation]{\textbf{Regulation (EU) 2023/1230}}. They are expressed here as derived component-level requirements for the \gls{acr:SUT} for verification purposes. They do not state that the Regulation applies directly to the component. The requirements are mapped to \hyperref[subsec:pren_50742]{\textbf{prEN 50742}} and, subsequently, to \hyperref[subsec:iec_62443]{\textbf{\gls{acr:IEC} 62443-4-2}} \glspl{acr:CR}, \glspl{acr:RE}, and component-type-specific instantiations such as \glspl{acr:EDR}. The three tables correspond to the three derivation stages and share a common requirement identifier (\gls{acr:RQ}-XXX) to preserve traceability across all stages.

\section{Security Requirements Derived from Regulation (EU) 2023/1230}
\label{app:sec:mvo_requirements}

Table \ref{tab:mvo-requirements} lists the 14 derived component-level requirements originating from Annex III, Sections 1.1.9 and 1.2.1 of Regulation (EU) 2023/1230. Each entry records the originating source clause, the requirement statement, the assigned Security Objective, the Requirement Type, and a rationale explaining the derivation from the regulatory text.



\renewcommand{\arraystretch}{1.25}
\begin{longtable}{|>{\bfseries}l|p{11cm}|}
	\caption{Security requirements derived from Regulation (EU) 2023/1230} 
	\label{tab:mvo-requirements} \\
	\hline
	\textbf{Field} & \textbf{Description} \\
	\hline
	\endfirsthead
	
	\multicolumn{2}{c}{\tablename\ \thetable{} continued from previous page} \\
	\hline
	\textbf{Field} & \textbf{Description} \\
	\hline
	\endhead
	
	
	\hline
	\endlastfoot
	
	\multicolumn{2}{|l|}{\textbf{Requirement RQ-001}} \\
	\hline
	Source & Annex III, Section 1.1.9 \\
	\hline
	Requirement & The system shall be designed so that the connection of another device, including any remote device that communicates with it, does not lead to a hazardous situation. \\
	\hline
	Security Objective & Integrity \\
	\hline
	\hline
	
	\multicolumn{2}{|l|}{\textbf{Requirement RQ-002}} \\
	\hline
	Source & Annex III, Section 1.1.9 \\
	\hline
	Requirement & The system shall protect hardware components transmitting signal or data relevant to connection or access to safety-critical software against corruption, whether accidental or intentional. \\
	\hline
	Security Objective & Integrity \\
	\hline
	\hline
	
	\multicolumn{2}{|l|}{\textbf{Requirement RQ-003}} \\
	\hline
	Source & Annex III, Section 1.1.9 \\
	\hline
	Requirement & The system shall collect evidence of any legitimate or illegitimate intervention in hardware components relevant for connection or access to safety-critical software. \\
	\hline
	Security Objective & Accountability \\
	\hline
	\hline
	
	\multicolumn{2}{|l|}{\textbf{Requirement RQ-004}} \\
	\hline
	Source & Annex III, Section 1.1.9 \\
	\hline
	Requirement & The system shall identify the software and data that are critical for compliance with the essential health and safety requirements. \\
	\hline
	Security Objective & Integrity \\
	\hline
	\hline
	
	\multicolumn{2}{|l|}{\textbf{Requirement RQ-005}} \\
	\hline
	Source & Annex III, Section 1.1.9 \\
	\hline
	Requirement & The system shall protect identified safety-critical software and data against corruption, whether accidental or intentional. \\
	\hline
	Security Objective & Integrity \\
	\hline
	\hline
	
	\multicolumn{2}{|l|}{\textbf{Requirement RQ-006}} \\
	\hline
	Source & Annex III, Section 1.1.9 \\
	\hline
	Requirement & The system shall identify the software installed on it that is necessary for safe operation. \\
	\hline
	Security Objective & Integrity \\
	\hline
	\hline
	
	\multicolumn{2}{|l|}{\textbf{Requirement RQ-007}} \\
	\hline
	Source & Annex III, Section 1.1.9 \\
	\hline
	Requirement & The system shall be able to provide the identification information of safety-relevant installed software at all times in an easily accessible form. \\
	\hline
	Security Objective & Availability \\
	\hline
	\hline
	
	\multicolumn{2}{|l|}{\textbf{Requirement RQ-008}} \\
	\hline
	Source & Annex III, Section 1.1.9 \\
	\hline
	Requirement & The system shall collect evidence of any legitimate or illegitimate intervention in the software installed on it. \\
	\hline
	Security Objective & Accountability \\
	\hline
	\hline
	
	\multicolumn{2}{|l|}{\textbf{Requirement RQ-009}} \\
	\hline
	Source & Annex III, Section 1.1.9 \\
	\hline
	Requirement & The system shall collect evidence of any modification of the installed software or its configuration. \\
	\hline
	Security Objective & Accountability \\
	\hline
	\hline
	
	\multicolumn{2}{|l|}{\textbf{Requirement RQ-010}} \\
	\hline
	Source & Annex III, Section 1.2.1 \\
	\hline
	Requirement & The control system shall withstand reasonably foreseeable malicious attempts from third parties that could lead to a hazardous situation. \\
	\hline
	Security Objective & Availability \\
	\hline
	\hline
	
	\multicolumn{2}{|l|}{\textbf{Requirement RQ-011}} \\
	\hline
	Source & Annex III, Section 1.2.1 \\
	\hline
	Requirement & The system shall prevent modifications to safety-relevant settings or rules, including those generated during a learning phase, where such modifications could lead to a hazardous situation. \\
	\hline
	Security Objective & Integrity \\
	\hline
	\hline
	
	\multicolumn{2}{|l|}{\textbf{Requirement RQ-012}} \\
	\hline
	Source & Annex III, Section 1.2.1 \\
	\hline
	Requirement & The system shall enable a tracing log of data generated in relation to an intervention for five years after the machinery or related product has been placed on the market or put into service. \\
	\hline
	Security Objective & Accountability \\
	\hline
	\hline
	
	\multicolumn{2}{|l|}{\textbf{Requirement RQ-013}} \\
	\hline
	Source & Annex III, Section 1.2.1 \\
	\hline
	Requirement & The system shall enable a tracing log of the versions of safety software uploaded after the machinery or related product has been placed on the market or put into service, for five years after each upload. \\
	\hline
	Security Objective & Accountability \\
	\hline
	\hline
	
	\multicolumn{2}{|l|}{\textbf{Requirement RQ-014}} \\
	\hline
	Source & Annex III, Section 1.2.1 \\
	\hline
	Requirement & The system shall restrict access to the tracing log data exclusively to demonstrating conformity further to a reasoned request from a competent national authority. \\
	\hline
	Security Objective & Confidentiality \\
	\hline
	\hline
	
\end{longtable}

\section{Mapping to prEN 50742}
\label{app:sec:mapping_pren50742}

Table \ref{tab:pren50742-mapping} maps the 14 requirements of Table \ref{tab:mvo-requirements} onto the corresponding clauses of \hyperref[subsec:pren_50742]{\textbf{prEN 50742}} (Safety of machinery - Protection against corruption). The mapping types are used as follows: \textbf{Partial} indicates that the clause addresses the obligation but does not fully cover it; \textbf{Supporting} indicates that the clause contributes context or a prerequisite but is not the operative requirement; and \textbf{No mapping} indicates that the draft standard explicitly excludes or does not address the obligation.
\begin{longtable}{|>{\bfseries}p{1cm}|p{2.5cm}|p{9cm}|>{\centering\arraybackslash}p{2cm}|}
	\caption{Mapping of Machinery Regulation requirements to prEN 50742} 
	\label{tab:pren50742-mapping} \\
	\hline
	\textbf{\gls{acr:RQ} ID} & \textbf{prEN 50742 clause} & \textbf{prEN 50742 requirement/topic} & \textbf{Mapping type} \\
	\hline
	\endfirsthead
	
	\multicolumn{4}{c}{\tablename\ \thetable{} continued from previous page} \\
	\hline
	\textbf{\gls{acr:RQ} ID} & \textbf{prEN 50742 clause} & \textbf{prEN 50742 requirement/topic} & \textbf{Mapping type} \\
	\hline
	\endhead
	
	\hline
	\endlastfoot
	RQ-001 &
	\S4.2; \newline \S4.3 &
	Design/construction of machinery shall not permit hazardous situations arising from connections (physical, logical, indirect). &
	\S4.2: Supporting \newline\S4.3: Supporting \\
	\hline
	
	
	RQ-002 &
	\S1;  \newline \S3.5;  \newline \S7.1;  \newline \S7.2.1 &
	Scope explicitly includes hardware components/interfaces (incl. remote-device interfaces) that transmit signals/data affecting machine safety; \S7.1 requires identification and protection of accessible interfaces by compensating/countermeasures. &
	\S1, \S3.5, \S7.2.1: Supporting\newline\S7.1: Partial \\
	\hline
	
	RQ-003 &
	\S3.4;  \newline \S7.3.3;  \newline\S7.4.3.4.5 &
	``Evidence'' definition includes physical tamper evidence (broken seals, damaged interface covers, e.g. RJ45); \S7.3.3 requires physical evidence where digital evidence collection is not reasonably practicable; \S7.4.3.4.5 is the only operative requirement mandating detection of physical manipulation (\gls{acr:SRSL}0-3).  &
	\S3.4: Supporting\newline\S7.3.3, \S7.4.3.4.5: Partial \\
	\hline
	
	RQ-004 &
	\S3.12;  \newline\S7.5 &
	``Critical data'' defined as data whose manipulation can increase hazard; \S7.5 requires identification of software versions/configuration data relevant for protection against falsification. &
	\S3.12: Supporting\newline\S7.5: Partial \\
	\hline
	
	RQ-005 &
	\S3.5;  \newline \S7.1;  \newline\S7.4.3.4.1 &
	\gls{acr:SRSL}-graded software/information integrity requirements, from checksum verification at start-up (\gls{acr:SRSL}1) to cryptographic hash/HMAC/CMAC verification during operation (\gls{acr:SRSL}3). &
	\S3.5, \S7.1: Supporting\newline\S7.4.3.4.1: Partial \\
	\hline
	
	
	RQ-006 &
	\S7.5;  \newline\S9 &
	Identification of software versions/configuration data ``relevant for protection against falsification'' must be available; \S9 additionally requires that user information include the identification of software versions and configuration data relevant for safety functions. &
	\S7.5: Partial\newline\S9: Supporting \\
	\hline
	
	RQ-007 &
	\S7.5; \newline \S9 &
	Identification data ``must be available on request'', in readable form, via an integrated system or an external tool; \S9 requires user information to include instructions for accessing this data and its modifications. Note: prEN 50742 requires availability ``on request'' rather than ``at all times'' as mandated by the regulation, so the mapping remains partial. &
	\S7.5: Partial\newline\S9: Supporting \\
	\hline
	
	RQ-008 &
	\S3.2;  \newline \S7.3.1;  \newline \S7.3.2;  \newline \S7.3.3;  \newline\S7.3.4 &
	Enumerated intervention types with minimum recorded data; \S7.3.4 specifies the log must at least include the records of the last intervention of each type. &
	\S3.2: Supporting\newline\S7.3.1, \S7.3.2, \S7.3.3, \S7.3.4: Partial\\
	\hline
	
	RQ-009 &
	\S3.2; \newline \S3.3;  \newline \S7.3.1;  \newline \S7.3.2;   \newline\S7.3.3; \newline \S7.3.4 &
	``Modification'' is defined as a change caused by an ``intervention'', incl. software version, sensor calibration, configuration data); \S7.3.1 b)/c) require identification means (version number, \gls{acr:CRC}) for SRESW/SRASW updates; \S7.3.4 specifies the associated minimum retention of at least five years per intervention type. &
	\S3.2, \S3.3: Supporting\newline\S7.3.1, \S7.3.2, \S7.3.3, \S7.3.4: Partial \\
	\hline
	
	RQ-010 &
	\S4.3;  \newline \S7.1;  \newline \S7.2.1;  \newline \S7.4.2;  \newline\S7.4.3 &
	Process principles: eliminate vulnerabilities that could lead to hazardous situations; otherwise mitigate; otherwise inform the user via user information of compensating measures; \S7.1 requires identification and protection of accessible interfaces, \S7.2.1 requires risk-based selection of security measures, and \S7.4.2/\S7.4.3 require determination and maintenance of \gls{acr:SRSL}-graded security requirements, all contributing to withstanding malicious third-party attempts. &
	\S4.3, \S7.2.1, \S7.4.2: Supporting\newline\S7.1, \S7.4.3: Partial \\
	\hline
	
	RQ-011 &
	\S7.4.3.3.1;  \newline \S3.2;  \newline \S3.3 &
	\gls{acr:SRSL}-graded authorisation requirements for interventions (\gls{acr:SRSL}1/2: authorisation required; \gls{acr:SRSL}3: privileged/role-based authorisation); the ``intervention'' definition names safety routines, safety variables, and calibration data as examples. Note: prEN 50742 contains no provision addressing learning-phase or self-evolving machinery behaviour; Annex ZZ correlation table (Table ZZ.1) does not correlate 1.2.1, Abs. 2, Point d) with any clause &
	\S7.4.3.3.1: Partial\newline\S3.2, \S3.3: Supporting \\
	\hline
	
	RQ-012 &
	\S7.3.1;  \newline \S7.3.3;  \newline \S7.3.4 &
	Log must contain at least the evidence of the last intervention of each type; records retained at least five years; log active at the latest from the time the machinery is put into service. &
	\S7.3.1, \S7.3.3, \S7.3.4: Partial \\
	\hline
	
	RQ-013 &
	\S7.3.1 &
	Explicit exclusion: versions of SRASW/SRESW binaries need not be stored in the tracing log. &
	No mapping \\
	\hline
	
	RQ-014 &
	\S7.3.5;  \newline \S7.3.3;  \newline \S7.4.3.3.1 &
	Logs must be ``properly secured'' against tampering; deletion only via an authorised procedure; digital evidence must be readable and accessible. &
	\S7.3.5, \S7.3.3: Supporting\newline\S7.4.3.3.1: Partial \\
	\hline
	
	
\end{longtable}

Where the draft standard explicitly excludes an obligation, as for RQ-013, the derivation proceeds directly from the Regulation; the absence of a standard clause does not remove the underlying obligation.


\section{Mapping to \gls{acr:IEC} 62443-4-2}
\label{app:sec:mapping_iec62443}
%TODO zeilen abstand ist komisch
\normalsize
Table \ref{tab:mvo-pren-iec62443-mapping} extends the derivation chain from \hyperref[subsec:pren_50742]{\textbf{prEN 50742}} to \gls{acr:IEC} 62443-4-2. Where no \hyperref[subsec:pren_50742]{\textbf{prEN 50742}} clause is available, the mapping proceeds directly from \hyperref[subsec:eu_machinery_regulation]{\textbf{Regulation (EU) 2023/1230}}. The mapped \glspl{acr:CR} are to be satisfied at the target Component \gls{acr:SL-C} declared in Section \ref{subsec:derivation_mapping}; applicable \glspl{acr:RE} are included accordingly. \gls{acr:IEC} 62443-4-2 organises its \glspl{acr:CR} under the seven \glspl{acr:FR} introduced in Section \ref{subsec:iec_62443}.

\footnotesize
\begin{longtable}{|>{\bfseries}p{1.0cm}|p{0.9cm}|p{2.4cm}|p{2.4cm}|p{4.8cm}|p{1.8cm}|}
	\caption{Mapping of derived requirements to \gls{acr:IEC} 62443-4-2} 
	\label{tab:mvo-pren-iec62443-mapping} \\
	\hline
	\textbf{\gls{acr:RQ} ID} & \textbf{\gls{acr:FR}} & \textbf{\gls{acr:CR}/\gls{acr:RE}/\gls{acr:EDR}} & \textbf{\gls{acr:IEC} 62443-4-2 topic} & \textbf{Mapping rationale} & \textbf{Type} \\
	\hline
	\endfirsthead
	
	\multicolumn{6}{c}{\tablename\ \thetable{} continued from previous page} \\
	\hline
	\textbf{\gls{acr:RQ} ID} & \textbf{\gls{acr:FR}} & \textbf{\gls{acr:CR}/\gls{acr:RE}/\gls{acr:EDR}} & \textbf{\gls{acr:IEC} 62443-4-2 topic} & \textbf{Mapping rationale} & \textbf{Type} \\
	\hline
	\endhead
	
	
	\hline
	\endlastfoot
	
	RQ-001 &
	\gls{acr:FR}3 &
	\gls{acr:CR} 3.1 &
	Communication integrity &
	\gls{acr:CR} 3.1 requires the component to verify the integrity of information transmitted over a communication channel, the closest component-level control for connection-related corruption. &
	Partial \\
	\hline
	
	RQ-002 &
	\gls{acr:FR}3 &
	\gls{acr:CR} 3.1; \newline \gls{acr:EDR} 3.11 &
	Communication integrity; Physical tamper resistance and detection (embedded device) &
	\gls{acr:CR} 3.1 covers the integrity of the information transmitted over such an interface, while \gls{acr:EDR} 3.11 covers protection of the interface hardware itself against unauthorised physical access and corruption. &
	Partial \\
	\hline
	
	RQ-003 &
	\gls{acr:FR}2; \gls{acr:FR}3 &
	\gls{acr:CR} 2.8; \newline \gls{acr:CR} 2.12; \newline \gls{acr:EDR} 3.11 &
	Auditable events; Non-repudiation; Physical tamper resistance and detection (embedded device) &
	prEN \S7.3.3 makes physical evidence (\gls{acr:EDR} 3.11) the fallback used only where digital evidence collection (\gls{acr:CR} 2.8) is not practicable. \gls{acr:CR} 2.8 supplies the primary digital-evidence/event-generation mechanism; \gls{acr:CR} 2.12 supplies the attribution property implied by ``legitimate or illegitimate'' (the evidence must be undeniably traceable to the causing action) &
	Partial \\
	\hline
	
	RQ-004 &
	\gls{acr:FR}7 &
	\gls{acr:CR} 7.8 &
	Control system component inventory &
	\gls{acr:CR} 7.8's component-inventory capability is the closest \gls{acr:CR}-level asset-identification control. --> \gls{acr:IEC} 62443-4-1: SUM-1/SM-9/SM-10 &
	Supporting \\
	\hline
	
	RQ-005 &
	\gls{acr:FR}3 &
	\gls{acr:CR} 3.4; \newline \gls{acr:EDR} 3.4 &
	Software and information integrity; Integrity of the boot process (embedded device) &
	prEN's \gls{acr:SRSL}-graded integrity framework (\S7.4.3.4) covers general software/information integrity (\S7.4.3.4.1 $\rightarrow$ \gls{acr:CR} 3.4) and, as a distinct sub-case, boot-process integrity (\S7.4.3.4.2).  &
	Partial \\
	\hline
	
	RQ-006 &
	\gls{acr:FR}7 &
	\gls{acr:CR} 7.8 &
	Control system component inventory &
	\gls{acr:CR} 7.8's inventory capability is the closest \gls{acr:CR}-level match for tracking installed-software identity. --> \gls{acr:IEC} 62443-4-1: SUM-1/SM-9/SM-1&
	Supporting \\
	\hline
	
	RQ-007 &
	\gls{acr:FR}7 &
	\gls{acr:CR} 7.1, \gls{acr:CR} 7.2, \gls{acr:CR} 7.8 &
	Control system component inventory &
	\gls{acr:CR} 7.8 supports retrieval of inventory information but does not mandate continuous, unsolicited availability, \gls{acr:CR}7.1 (DoS protection), \gls{acr:CR}7.2 Resource Management --> \gls{acr:IEC} 62443-4-1: SUM-1/SM-9/SM-1 &
	Supporting \\
	\hline
	
	RQ-008 &
	\gls{acr:FR}2 &
	\gls{acr:CR} 2.8; \newline \gls{acr:CR} 2.12 &
	Auditable events; Non-repudiation &
	prEN \S7.3.1/\S7.3.2 anchor the enumerated intervention types and minimum recorded data to \gls{acr:CR} 2.8 (event generation); the ``legitimate or illegitimate'' attribution qualifier matches \gls{acr:CR} 2.12 (non-repudiation). &
	\gls{acr:CR} 2.8: Partial\newline \gls{acr:CR} 2.12: Supporting \\
	\hline
	
	RQ-009 &
	\gls{acr:FR}2; \gls{acr:FR}3 &
	\gls{acr:CR} 2.8; \newline \gls{acr:CR} 2.12; \newline \gls{acr:CR} 3.4 &
	Auditable events; Non-repudiation; Software and information integrity &
	prEN's ``modification'' definition and the version-number/\gls{acr:CRC} identification means combine an audit obligation (\gls{acr:CR} 2.8) with an integrity-verification obligation (\gls{acr:CR} 3.4). Because a modification is defined as a change caused by an intervention (\S3.2), the attribution qualifier extends to modification evidence, which \gls{acr:CR} 2.12 (non-repudiation) addresses. &
	\gls{acr:CR} 2.8, \gls{acr:CR} 3.4: Partial \newline \gls{acr:CR} 2.12: Supporting \\
	\hline
	
	RQ-010 &
	\gls{acr:FR}7; \newline (\gls{acr:FR}1; \newline \gls{acr:FR}2; \newline 
	\gls{acr:FR}3)&
	\gls{acr:CR} 7.1 (1) &
	Protection against denial-of-service events &
	The regulation's requirement that control systems withstand reasonably foreseeable malicious third-party attempts corresponds, at the component level, to \gls{acr:CR} 7.1 (protection against denial-of-service events) together with its \gls{acr:RE}(1) (mitigation of communication-load flooding). &
	Partial \\
	\hline
	
	RQ-011 &
	\gls{acr:FR}1; \gls{acr:FR}2; \gls{acr:FR}3 &
	\gls{acr:CR} 1.1; \newline \gls{acr:CR} 1.2; \newline \gls{acr:CR} 2.1 (2); \newline \gls{acr:CR} 3.4 &
	Human/device identification and authentication; Authorisation enforcement (role-based); Software and information integrity &
	prEN's \gls{acr:SRSL}-graded authorisation requirement for interventions corresponds to \gls{acr:CR} 2.1, with role/task-based authorisation at \gls{acr:SRSL}3 (\S7.4.3.3.1) matching \gls{acr:CR} 2.1's \gls{acr:RE}(2) (mapping of permissions to roles). \gls{acr:CR} 1.1 addresses authentication of the human-operator path, \gls{acr:CR} 1.2 (identification and authentication of software processes and devices) additionally applies to that path. \gls{acr:CR} 3.4 is included because ``preventing modifications to safety-relevant settings or rules'' is fundamentally a configuration-integrity property. The learning-phase/self-evolving-behaviour aspect has no counterpart in \gls{acr:IEC} 62443-4-2. &
	\gls{acr:CR} 1.1, \gls{acr:CR} 1.2, \gls{acr:CR} 3.4: Supporting \newline \gls{acr:CR} 2.1 (2): Partial\\
	\hline
	
	RQ-012 &
	\gls{acr:FR}2 &
	\gls{acr:CR} 2.8; \newline \gls{acr:CR} 2.9 &
	Auditable events; Audit storage capacity &
	prEN's requirement that the log retain at least the evidence of the last intervention of each type combines the event-generation obligation of \gls{acr:CR} 2.8 with the storage-capacity obligation of \gls{acr:CR} 2.9. \gls{acr:IEC} 62443-4-2 defines no Component Requirement mandating a specific retention duration. &
	\gls{acr:CR} 2.8, \gls{acr:CR} 2.9: Partial \\
	\hline
	
	RQ-013 &
	\gls{acr:FR}2 &
	\gls{acr:CR} 2.8; \newline \gls{acr:CR} 2.9 &
	Auditable events (configuration changes); Audit storage capacity &
	\gls{acr:CR} 2.8's auditable-event categories include configuration changes and automation-system events, which are broad enough to cover the upload of a new safety-software version as a recordable event, and \gls{acr:CR} 2.9 provides the associated storage-capacity capability (without mandating the specific five-year, per-upload retention period). &
	\gls{acr:CR} 2.8, \gls{acr:CR} 2.9: Partial \\
	\hline
	
	RQ-014 &
	\gls{acr:FR}2; \gls{acr:FR}3; \gls{acr:FR}4; \gls{acr:FR}6 &
	\gls{acr:CR} 2.1; \newline \gls{acr:CR} 3.9; \newline \gls{acr:CR} 4.1; \newline \gls{acr:CR} 6.1 (supporting) &
	Authorisation enforcement; Protection of audit information; Data confidentiality; Audit log accessibility (supporting) &
	\gls{acr:CR} 2.1 (authorisation enforcement) is the primary control restricting who may access the log and for what purpose, and \gls{acr:CR} 4.1 (confidentiality of information at rest for which read authorisation is explicitly required). \gls{acr:CR} 3.9 additionally protects the log against tampering and unauthorised deletion. \gls{acr:CR} 6.1 (ability to access event logs) is included only as a supporting control, since it ensures that an already-authorised party can retrieve the log on request. &
	\gls{acr:CR} 2.1: Partial\newline \gls{acr:CR} 3.9, \gls{acr:CR} 4.1, \gls{acr:CR} 6.1: Supporting \\
	\hline
	
\end{longtable}
\normalsize

\chapter{Secure Development Lifecycle}
\label{app:SecureDevelopmentLifecycle}


Figure \ref{fig:secure-system-development-lifecycle-main-overview} summarises a high-level reference model in which the activities shown are decomposed into further sub-activities. 
The model is based on the process-oriented security practices of \gls{acr:IEC} 62443-4-1 \cite{DIN_EN_IEC_62443_4_1_2018} and on complementary lifecycle-oriented security approaches \cite{NiemannPuls2023}. 

The lifecycle demonstrates that cybersecurity is not limited to a single verification activity prior to product release. Instead, security is addressed throughout the development and operational lifecycle \cite[pp. 5]{NiemannPuls2023} \cite[pp. 68]{Wurm2022}. The process begins with the definition of the scope and security objectives. A \hyperref[subsec:tara]{\textbf{TARA}}, is performed to identify and assess relevant cybersecurity risks and to determine the required security level. Based on the assessment results, detailed countermeasures are derived and incorporated into the security architecture. The resulting architectural constraints provide input to the development phase, during which the specified security measures are implemented \cite[pp. 8]{NiemannPuls2023}. 

The implemented security measures are evaluated by means of static and dynamic analysis. The verification phase further comprises the definition and preparation of the security test strategy, the development of test cases, and the execution of the corresponding security tests. Identified findings are returned to the development phase for remediation and subsequent re-evaluation. The verification results provide the evidence required to support the release decision \cite[pp. 788]{Muller2026} \cite[pp. 6, 16]{NiemannPuls2023}. Following release, continuous monitoring is applied to identify newly discovered vulnerabilities, emerging threats, and security-relevant incidents. The resulting observations may initiate a reassessment or an incident-response process \cite[pp. 18]{NiemannPuls2023}.

The dashed connections shown in Figure \ref{fig:secure-system-development-lifecycle-main-overview} represent feedback paths and emphasise that the lifecycle is iterative rather than strictly sequential. Findings identified during security testing are returned to secure implementation for remediation. Architectural constraints may require the assumptions or results of the risk assessment to be reconsidered. In addition, information obtained through continuous monitoring may trigger a reassessment or incident response, thereby re-entering the lifecycle at the risk-assessment stage. These feedback paths support the continuous refinement of the risk assessment, security architecture, implementation, and verification activities 


\begin{figure}[H]
	\centering
	\includegraphics[width=\linewidth]{"images/Secure System Development Lifecycle-Main Overview.drawio"}
	\caption[High-level overview of the secure development lifecycle]{High-level overview of the secure development lifecycle (Own illustration)}
	\label{fig:secure-system-development-lifecycle-main-overview}
	\footnotetext{Author's own illustration based on \cite{DIN_EN_IEC_62443_4_1_2018} and \cite{NiemannPuls2023}}
\end{figure}

\paragraph{\gls{acr:RACI} Governance of the Verification Lifecycle}
\label{sec:raci_governance}

To ensure strict organisational accountability, responsibility for the testing phases is governed through a \gls{acr:RACI} matrix. This management method is described as ensuring exactly one role is accountable for a given activity \cite{cybersierra_security_raci_2026, su13041806}. In this scheme, the responsible role performs the activity, the accountable role owns and approves the outcome, consulted roles provide input before a decision, and informed roles are notified of the result. Six generic roles are distinguished for this purpose \cite[p. 6]{Schmidt2026CybersecurityTestProcess}. 

\begin{itemize}
	\item \textbf{Security Engineering:} Holds technical accountability for test scope and execution integrity.
	\item \textbf{System/Requirements Engineering:} Consulted where a requirement or asset boundary is affected.
	\item \textbf{Test Engineering:} Responsible for procedure preparation, execution, and verdict assignment.
	\item \textbf{Software/Firmware Development:} Consulted on findings and tasked with implementing corrective actions.
	\item \textbf{Quality/Compliance Management:} Responsible for coverage review and evidence governance.
	\item \textbf{Release Authority:} Accountable for the final release decision.
\end{itemize}

Table \ref{tab:raci_matrix_verification} defines the distribution of these responsibilities across the verification phases.

\begin{table}[H]
	\caption[\gls{acr:RACI} matrix for the verification lifecycle]{\gls{acr:RACI} matrix for the verification lifecycle (Adapted from \cite{cybersierra_security_raci_2026})}
	\label{tab:raci_matrix_verification}
	\centering
	\begin{tabular*}{\textwidth}{@{\extracolsep{\fill}} l c c c c c c @{}}
		\toprule
		\textbf{Activity} & \makecell{\textbf{Sec.}\\\textbf{Eng.}} & \makecell{\textbf{Sys/Req}\\\textbf{Eng.}} & \makecell{\textbf{Test}\\\textbf{Eng.}} & \makecell{\textbf{SW/FW}\\\textbf{Dev.}} & \makecell{\textbf{QA/}\\\textbf{Compl.}} & \makecell{\textbf{Release}\\\textbf{Auth.}} \\ \midrule
		Design & A & C & R & I & I & I \\ \addlinespace
		Preparation & C & I & A/R & C & I & I \\ \addlinespace
		Execution & A & I & R & C & I & I \\ \addlinespace
		Evidence/Eval. & C & I & R & C & A & I \\ \addlinespace
		Release Decision & I & I & I & I & C & A/R \\
		\bottomrule
	\end{tabular*}
\end{table}

As illustrated in Table \ref{tab:raci_matrix_verification}, accountability shifts between phases. It follows the organisational ownership of a phase's primary deliverable rather than a fixed functional assignment \cite{cybersierra_security_raci_2026}. Distinguishing the role responsible for executing a verification activity from the role accountable for closing it out reflects the practical division of labour documented in secure software engineering \cite[p. 5]{219394}. Ultimately, these role assignments directly support \gls{acr:IEC} 62443-4-1 Practice \gls{acr:SVV}. This standard mandates a rigorous process demonstrating that the product's security functions meet all specified requirements \cite[pp. 37]{DIN_EN_IEC_62443_4_1_2018}.

\newpage


\chapter{Tool Commands}
\label{app:tool_commands}
This appendix lists the tool commands used for the activities executed under the \gls{acr:NV1} network position, as reported in Section \ref{sec:executed_tests}. The placeholder \texttt{\$SRIO} denotes the \gls{acr:IP} address of the \gls{acr:SRIO}, \texttt{192.168.0.2}.
All commands were executed from the Kali Linux test system under \gls{acr:NV1} in the isolated laboratory environment described in Section \ref{subsec:realized_environment}. The commands are reproduced as run, and the file names below are the names of the stored evidence artefacts produced during the test campaign.
\begin{verbatim}
	export SRIO=192.168.0.2
\end{verbatim}
The output file names correspond to the stored evidence.


\section{Attack-Surface Enumeration}
\label{app:cmd_attacksurface}
The scan was executed on two separate dates to confirm reproducibility, both with the same result (Section \ref{subsec:res_attack_surface}). The commands and file names of the later run are given below; the earlier run is stored under the \texttt{cm001\_srio\_*} basename with otherwise identical commands.
\begin{verbatim}
	
	# Full TCP scan with service detection
	sudo nmap -sS -sV -p- -T3 -oA cm_rq001-02_tcp_srio $SRIO
	# UDP top-ports scan
	sudo nmap -sU --top-ports 200 -oA cm_rq001-02 $SRIO
	# Web-server scan
	nikto -h http://$SRIO -output cm_rq001-02_srio.html -Format html
\end{verbatim}

\section{Identification-Surface Inventory}
\label{app:cmd_inventory}
\begin{verbatim}
	for p in deviceinfo/hwversion deviceinfo/hwrevision deviceinfo/swrevision \
	deviceinfo/swinfo/cpubootloaderversion deviceinfo/swinfo/scpuversion \
	deviceinfo/swinfo/scpubootloaderversion firmware/version \
	fieldbussetup/fieldbusfirmware; do
	echo "== $p =="
	curl -s http://$SRIO/$p/getdata
	echo
	done | tee cm_rq004-01_inventory.txt
\end{verbatim}

\section{Cleartext Transmission of the Identification Data}
\label{app:cmd_cleartext}

\begin{verbatim}
	sslscan --show-certificate $SRIO:80 | tee cm_rq006-03_sslscan.txt
\end{verbatim}

\section{Intervention Evidence and Time Attribution}
\label{app:cmd_evidence}
\begin{verbatim}
	# Read the error log and the system tick counter before and after each event
	curl -s http://$SRIO/devicestatus/errorlog/loglist | tee errorlog.json
	curl -s http://$SRIO/systemtime/systick/getdata  | tee systick.txt
	# Manual sequence: baseline; short-circuit pins 1 and 3 of a digital input;
	# remove power for a cold restart; 
	induce the second short-circuit; 
	repeat the reads.
\end{verbatim}

\section{Availability under Connection Load}
\label{app:cmd_availability}

\begin{verbatim}
	# Application-layer load test at 50 and at 3 concurrent connections
	siege -c 50 -t 1M -v http://$SRIO/deviceinfo/getdata \
	--log=cmrq010_siege.log
	siege -c 3  -t 1M -v http://$SRIO/deviceinfo/getdata \
	--log=cmrq010_3connections_siege.log
	# Connection-hold procedure: occupy the available HTTP slots, then issue
	bash
	export SRIO=192.168.0.2
	echo " Baseline: Connection working normally" 
	| tee cm_rq007-03_00_baseline.txt
	( time curl -s --max-time 5 http://$SRIO/deviceinfo/swrevision/getdata )
	>> cm_rq007-03_00_baseline.txt 2>&1
	cat cm_rq007-03_00_baseline.txt
	
	echo "Block both slots"
	for i in 1 2; do
	(exec 3<>/dev/tcp/$SRIO/80; printf 'GET /deviceinfo HTTP/1.1\r\nHost:
	x\r\n\r\n' >&3; sleep 600) &
	done
	jobs | tee cm_rq007-03_01_jobs.txt
	
	echo "3rd attempt while both slots are occupied (expected: timeout)" 
	| tee cm_rq007-03_02_blocked.txt
	( time curl -s --max-time 5 http://$SRIO/deviceinfo/swrevision/getdata ) 
	>> cm_rq007-03_02_blocked.txt 2>&1
	cat cm_rq007-03_02_blocked.txt
	
	echo "Free up one slot"
	kill %1
	echo "Retry after freeing up slot (expected: successful again)" 
	| tee cm_rq007-03_03_after_release.txt
	( time curl -s --max-time 5 http://$SRIO/deviceinfo/swrevision/getdata ) 
	>> cm_rq007-03_03_after_release.txt 2>&1
	cat cm_rq007-03_03_after_release.txt
\end{verbatim}

\section{Layer-Two and Network Flooding}
\label{app:cmd_flood}
\begin{verbatim}
	sudo hping3 --flood --rand-source $SRIO
	sudo macof -i eth1
\end{verbatim}

\section{Robustness under Malformed Input}
\label{app:cmd_robustness}
\begin{verbatim}
	# boofuzz fuzzing script 
	# File: cm010_fuzz.py
	from boofuzz import *
	
	session = Session(target=Target(connection=TCPSocketConnection("$SRIO", 80)))
	s_initialize("request")
	s_string("USER")        # field to fuzz
	s_delim(" ")
	s_string("value")
	session.connect(s_get("request"))
	session.fuzz()
	
	# Run it:
	python3 cm010_fuzz.py
	
\end{verbatim}

\section{Access Control and Configuration Interfaces}
\label{app:cmd_access}
\begin{verbatim}
	# Write probe against safety-relevant IoT-Core nodes
	curl -s -X POST http://$SRIO/safecom/f_address/setdata \
	-d '{"value":123}' -w "%{http_code}\n"
	curl -s -X POST http://$SRIO/io/do/port7/pin4/digital_output/setdata \
	-d '{"value":true}' -w "%{http_code}\n"
	curl -s -X POST http://$SRIO/devicetag/applicationtag/setdata \
	-d '"factoryXYZ"' -w "%{http_code}\n"
	# Active web scan through a local proxy
	# Disabling the API key is acceptable only in the isolated laboratory.
	zaproxy -daemon -host 127.0.0.1 -port 8090 -config api.disablekey=true
	curl "http://127.0.0.1:8090/JSON/spider/action/scan/?url=http://$SRIO/"
\end{verbatim}
